\section{Distributed Ledger Technology dan IOTA}\label{sec:dlt-iota}

\subsection{Distributed Ledger Technology}

\subsubsection{Definisi dan Konsep Dasar}

Pada subbab sebelumnya, telah dibahas bahwa rekam medis elektronik (RME) konvensional umumnya dibangun dengan arsitektur terpusat yang menyebabkan risiko \emph{single point of failure}, keterbatasan interoperabilitas lintas institusi, serta tantangan integritas dan auditabilitas data \parencite{Ettaloui_2024_blockchain_ehr_slr,Mayer_2020_ehr_blockchain_slr}. 
Salah satu pendekatan yang banyak dikaji untuk mengatasi keterbatasan tersebut adalah \emph{Distributed Ledger Technology} (DLT), yaitu keluarga teknologi pencatatan transaksi terdistribusi yang dipelihara bersama oleh sejumlah node dalam sebuah jaringan tanpa otoritas tunggal \parencite{Mayer_2020_ehr_blockchain_slr}. 
Blockchain merupakan bentuk DLT yang paling banyak diaplikasikan, di mana transaksi dikelompokkan ke dalam blok yang saling terhubung secara kriptografis sehingga membentuk rantai blok (\emph{chain}) yang bersifat \emph{append-only} dan sulit untuk dimodifikasi tanpa terdeteksi \parencite{Ettaloui_2024_blockchain_ehr_slr,Mayer_2020_ehr_blockchain_slr}. 

Dalam konteks kesehatan, sejumlah tinjauan sistematis mencatat bahwa DLT digunakan terutama sebagai \emph{trust layer} untuk penyimpanan hash, metadata, dan jejak audit akses data kesehatan, sementara data klinis berukuran besar seperti citra medis tetap diletakkan di penyimpanan di luar rantai (\emph{off-chain}) \parencite{Ettaloui_2024_blockchain_ehr_slr,Mayer_2020_ehr_blockchain_slr}. 
Dengan demikian, DLT tidak dimaksudkan menggantikan seluruh fungsi RME, melainkan menyediakan lapisan integritas, non-\emph{repudiation}, dan koordinasi akses di atas sistem-sistem yang sudah ada \parencite{Yang_2021_blockchain_architecture_ehr,Ettaloui_2024_blockchain_ehr_slr}. 

\subsubsection{Karakteristik Utama DLT pada Domain Kesehatan}

Literatur mengenai EHR berbasis blockchain mengidentifikasi sejumlah karakteristik kunci DLT yang relevan untuk pengelolaan data kesehatan \parencite{Ettaloui_2024_blockchain_ehr_slr,Mayer_2020_ehr_blockchain_slr}. 
Pertama, sifat desentralisasi dan distribusi ledger di banyak node mengurangi ketergantungan pada satu pihak penyedia, sehingga menurunkan risiko kegagalan tunggal dan memfasilitasi skenario pertukaran data antar penyedia layanan kesehatan \parencite{Mayer_2020_ehr_blockchain_slr}. 
Kedua, penggunaan kriptografi kunci publik dan tanda tangan digital memungkinkan otentikasi, otorisasi, serta penandaan identitas entitas (pasien, tenaga kesehatan, dan institusi) secara kriptografis, yang penting untuk menerapkan \emph{fine-grained access control} dan pelacakan siapa melakukan apa terhadap data \parencite{Ettaloui_2024_blockchain_ehr_slr,Yang_2021_blockchain_architecture_ehr}. 

Karakteristik lain yang menonjol adalah imutabilitas, yaitu sifat ledger yang tidak dapat diubah tanpa konsensus, sehingga setiap perubahan atau akses terhadap data meninggalkan jejak permanen yang dapat diaudit \parencite{Ettaloui_2024_blockchain_ehr_slr,Mayer_2020_ehr_blockchain_slr}. 
Dengan cara ini, DLT dapat digunakan sebagai sumber kebenaran (\emph{source of truth}) atas riwayat operasi terhadap rekam medis, termasuk pembuatan, pembaruan, dan akses data oleh berbagai pihak \parencite{Yang_2021_blockchain_architecture_ehr}. 
Tinjauan sistematis juga menekankan peran DLT dalam meningkatkan transparansi dan akuntabilitas, sembari tetap mendukung privasi melalui teknik seperti enkripsi, pseudonimisasi, hingga model penyimpanan hybrid on-chain/off-chain \parencite{Mayer_2020_ehr_blockchain_slr,Ettaloui_2024_blockchain_ehr_slr}. 

Sejalan dengan itu, berbagai arsitektur EHR berbasis blockchain mengusulkan penggunaan \emph{smart contract} sebagai komponen logika di atas ledger untuk merepresentasikan kebijakan akses, persetujuan pasien, dan aturan penggunaan data \parencite{Yang_2021_blockchain_architecture_ehr,Rawat2025_Blockchain_Secure_EHR}. 
\emph{Smart contract} tersebut dieksekusi secara deterministik pada node-node di jaringan DLT sehingga dapat mengotomatiskan verifikasi hak akses, pencatatan aktivitas, dan penerapan kebijakan tanpa campur tangan manual yang rentan kesalahan \parencite{Yang_2021_blockchain_architecture_ehr}. 

\subsubsection{Cara Kerja Distributed Ledger Technology pada Rekam Medis Elektronik}

Secara teknis, blockchain sebagai salah satu implementasi DLT bekerja dengan merekam transaksi ke dalam blok yang saling terhubung secara kriptografis sehingga membentuk rantai blok (\emph{blockchain}) yang bersifat \emph{append-only} \parencite{Ettaloui_2024_blockchain_ehr_slr,Mayer_2020_ehr_blockchain_slr}. 
Setiap blok terdiri atas bagian header dan body, di mana header umumnya memuat hash blok sebelumnya, timestamp, nilai \emph{nonce}, serta Merkle root atas sekumpulan transaksi di dalam blok, sedangkan body berisi daftar transaksi atau peristiwa yang ingin dicatat pada ledger \parencite{Yang_2021_blockchain_architecture_ehr}. 
Penggunaan hash blok sebelumnya pada header menyebabkan setiap blok bergantung pada blok sebelumnya, sehingga perubahan terhadap isi blok akan mengubah hash dan merusak seluruh rantai setelahnya, yang menjadikan usaha pemalsuan riwayat transaksi menjadi sangat sulit tanpa menguasai mayoritas kekuatan komputasi jaringan \parencite{Mayer_2020_ehr_blockchain_slr}. 

Untuk memastikan semua node di jaringan memiliki pandangan yang konsisten terhadap ledger, blockchain menggunakan mekanisme konsensus terdistribusi \parencite{Ettaloui_2024_blockchain_ehr_slr}. 
Pada blockchain publik, \emph{Proof of Work} (PoW) merupakan salah satu mekanisme konsensus yang paling awal, di mana penambang bersaing untuk menemukan nilai \emph{nonce} yang menghasilkan hash blok dengan karakteristik tertentu (misalnya sejumlah bit nol di awal), dan blok pertama yang memenuhi kriteria tersebut diterima sebagai blok berikutnya setelah diverifikasi oleh node lain \parencite{Yang_2021_blockchain_architecture_ehr,Mayer_2020_ehr_blockchain_slr}. 
Pendekatan ini memberikan keamanan yang kuat namun menuntut konsumsi energi dan sumber daya komputasi yang tinggi, serta membatasi throughput karena tingkat produksi blok ditahan pada laju tertentu untuk menjaga stabilitas jaringan \parencite{Ettaloui_2024_blockchain_ehr_slr}. 

Pada konteks EHR, banyak rancangan lebih memilih blockchain berizin (\emph{permissioned blockchain}) yang menggunakan mekanisme konsensus alternatif seperti \emph{Practical Byzantine Fault Tolerance} (PBFT) atau variasinya, di mana node yang dipercaya (misalnya institusi kesehatan) berpartisipasi dalam protokol voting untuk menyepakati blok baru \parencite{Mayer_2020_ehr_blockchain_slr,Yang_2021_blockchain_architecture_ehr}. 
Pendekatan ini dapat meningkatkan throughput dan menurunkan latensi dibandingkan PoW dengan mengorbankan sebagian tingkat desentralisasi, namun secara praktis lebih sesuai untuk lingkungan kesehatan di mana partisipan jaringan dapat diidentifikasi dan diatur \parencite{Ettaloui_2024_blockchain_ehr_slr}. 

Dalam berbagai arsitektur EHR berbasis blockchain, data medis berukuran besar umumnya tidak disimpan langsung di ledger, melainkan pada penyimpanan off-chain seperti basis data relasional, penyimpanan dokumen, atau sistem berkas terdistribusi \parencite{Ettaloui_2024_blockchain_ehr_slr,Mayer_2020_ehr_blockchain_slr}. 
Ledger menyimpan hash atau pointer terhadap berkas EHR, beserta metadata seperti identitas kriptografis pemilik data, pihak yang mengakses, waktu akses, dan jenis operasi yang dilakukan, sehingga ledger bertindak sebagai sumber kebenaran atas riwayat peristiwa tanpa mengekspos isi data sensitif \parencite{Yang_2021_blockchain_architecture_ehr,Rawat2025_Blockchain_Secure_EHR}. 
\emph{Smart contract} dimanfaatkan untuk mengimplementasikan logika kontrol akses, misalnya dengan memverifikasi apakah alamat atau sertifikat tertentu berhak melakukan operasi baca atau tulis, dan secara otomatis menambahkan catatan audit baru pada ledger ketika operasi tersebut dilakukan \parencite{Yang_2021_blockchain_architecture_ehr,Rawat2025_Blockchain_Secure_EHR}. 
Dengan demikian, kombinasi penyimpanan on-chain dan off-chain memungkinkan pemisahan yang jelas antara lapisan penyimpanan klinis dan lapisan integritas serta audit yang disediakan oleh DLT \parencite{Ettaloui_2024_blockchain_ehr_slr}. 

\subsubsection{Tantangan DLT dan Blockchain Konvensional untuk EHR}

Meskipun menawarkan berbagai keunggulan, pemanfaatan blockchain konvensional pada EHR menimbulkan tantangan signifikan yang perlu diperhatikan \parencite{Mayer_2020_ehr_blockchain_slr,Ettaloui_2024_blockchain_ehr_slr}. 
Dari sisi kinerja, beberapa studi melaporkan keterbatasan throughput dan meningkatnya latensi konfirmasi ketika jumlah transaksi dan ukuran data bertambah, khususnya pada blockchain publik dengan mekanisme konsensus intensif seperti \emph{Proof of Work} \parencite{Mayer_2020_ehr_blockchain_slr}. 
Keterbatasan ini berpotensi menjadi penghalang ketika sistem harus menangani aliran data frekuensi tinggi, misalnya dari perangkat IoT kesehatan atau sumber PGHD yang kontinu \parencite{Ettaloui_2024_blockchain_ehr_slr}. 

Selain itu, integrasi blockchain dengan RME memerlukan kepatuhan terhadap regulasi perlindungan data seperti GDPR dan HIPAA, termasuk prinsip \emph{right to be forgotten} yang pada praktiknya sulit dipenuhi pada ledger yang imutabel \parencite{Ettaloui_2024_blockchain_ehr_slr,Mayer_2020_ehr_blockchain_slr}. 
Pendekatan yang umum diadopsi adalah menyimpan data medis sensitif di luar rantai dan hanya menempatkan hash atau pointer pada ledger, namun desain ini menuntut tata kelola penyimpanan off-chain yang kuat agar tidak menimbulkan titik kelemahan baru \parencite{Yang_2021_blockchain_architecture_ehr}. 
Tinjauan sistematis juga mencatat tantangan interoperabilitas, karena sebagian besar sistem EHR yang ada menggunakan standar yang beragam sehingga memerlukan pemetaan ke format bersama (misalnya FHIR) sebelum dapat dipadukan dalam arsitektur DLT \parencite{Mayer_2020_ehr_blockchain_slr,Ettaloui_2024_blockchain_ehr_slr}. 

Aspek lain yang tidak kalah penting adalah biaya dan kompleksitas operasional \parencite{Ettaloui_2024_blockchain_ehr_slr,Rawat2025_Blockchain_Secure_EHR}. 
Implementasi dan pemeliharaan jaringan blockchain, terutama jika melibatkan banyak institusi, membutuhkan infrastruktur komputasi dan penyimpanan yang tidak sedikit, serta mekanisme konsensus yang disepakati bersama antara pemangku kepentingan kesehatan \parencite{Yang_2021_blockchain_architecture_ehr}. 
Sebagian penelitian mengusulkan penggunaan blockchain berizin (\emph{permissioned blockchain}) dan teknik optimasi seperti sharding atau off-chain processing untuk mengurangi beban, namun pendekatan ini biasanya mengorbankan tingkat desentralisasi tertentu \parencite{Mayer_2020_ehr_blockchain_slr}. 

\subsubsection{Ringkasan Peran DLT dalam Konteks Tesis}

Dari uraian di atas dapat disimpulkan bahwa DLT, khususnya blockchain, menyediakan fondasi konseptual untuk membangun lapisan kepercayaan (\emph{trust layer}) di atas RME dengan menyediakan integritas, auditabilitas, dan mekanisme kontrol akses yang lebih kuat \parencite{Ettaloui_2024_blockchain_ehr_slr,Mayer_2020_ehr_blockchain_slr}. 
Namun, keterbatasan skalabilitas, biaya, dan isu regulasi pada blockchain konvensional menunjukkan perlunya eksplorasi platform DLT alternatif yang lebih selaras dengan karakteristik data kesehatan modern, terutama yang bersumber dari perangkat IoT dan PGHD \parencite{Ettaloui_2024_blockchain_ehr_slr}. 
Kebutuhan inilah yang kemudian mendorong pemilihan IOTA, sebagai salah satu varian DLT non-blockchain, untuk menjadi basis ledger dalam kerangka DecMed yang akan dibahas pada subbab berikutnya.

\subsection{IOTA dan Tangle}

\subsubsection{Definisi dan Konsep Dasar IOTA}

IOTA adalah platform DLT yang dirancang untuk mendukung skenario Internet of Things (IoT) dengan memanfaatkan struktur data \emph{Directed Acyclic Graph} (DAG) yang disebut Tangle, bukan rantai blok linier seperti pada blockchain \parencite{Elevating_2025_iota_ehealth,FromTangleToRebased_2024}. 
Alih-alih mengandalkan penambang khusus, setiap transaksi baru pada Tangle diwajibkan untuk melakukan verifikasi terhadap sejumlah transaksi sebelumnya dan menjalankan perhitungan kecil sebagai bentuk bukti kerja ringan, sehingga mekanisme konsensus bersifat terdistribusi di antara partisipan jaringan \parencite{FromTangleToRebased_2024}. 
Desain ini memungkinkan IOTA untuk menghilangkan biaya transaksi (\emph{feeless}) dan secara teoretis meningkatkan throughput ketika tingkat partisipasi jaringan bertambah, yang merupakan karakteristik penting bagi skenario IoT berdaya dan sumber daya terbatas \parencite{Elevating_2025_iota_ehealth,FromTangleToRebased_2024}. 

Seiring perkembangan, IOTA mengusulkan Tangle 2.0 dan arsitektur Rebased yang bertujuan meningkatkan desentralisasi dan keamanan dengan mendesain ulang komponen inti seperti mekanisme pemilihan tip dan cara jaringan mencapai konsensus \parencite{FromTangleToRebased_2024}. 
Analisis komparatif menunjukkan bahwa arsitektur baru tersebut memperbaiki sejumlah kelemahan desain generasi awal, seperti potensi sentralisasi dan kerentanan terhadap serangan tertentu, sehingga meningkatkan kesesuaian IOTA untuk aplikasi IoT kritikal \parencite{FromTangleToRebased_2024}. 

\subsubsection{Cara Kerja IOTA dan Tangle}

Berbeda dengan blockchain yang mengelompokkan transaksi ke dalam blok, IOTA menggunakan struktur grafik asiklik terarah (\emph{Directed Acyclic Graph}, DAG) yang disebut Tangle, di mana setiap transaksi direpresentasikan sebagai simpul (node) dalam graf dan saling mereferensikan transaksi lain \parencite{TheTangle_Popov_2017,FromTangleToRebased_2024}. 
Secara garis besar, ketika sebuah node ingin menerbitkan transaksi baru, ia harus memilih sejumlah transaksi yang belum terkonfirmasi penuh (\emph{tips}) dan memverifikasi validitasnya sebelum transaksi baru tersebut boleh ditambahkan ke Tangle \parencite{TheTangle_Popov_2017}. 
Pemilihan tips pada generasi awal Tangle biasanya dilakukan dengan algoritme berbasis rantai Markov (\emph{Markov Chain Monte Carlo}, MCMC) yang menelusuri graf dari transaksi terkini ke masa lalu dan memilih tips berdasarkan bobot akumulatif agar transaksi yang konsisten dan banyak didukung lebih sering dipilih \parencite{TheTangle_Popov_2017}. 

Setelah tips dipilih, node memeriksa keabsahan transaksi yang direferensikan, misalnya memastikan tanda tangan digital benar dan tidak terjadi konflik saldo atau \emph{double-spending} \parencite{TheTangle_Popov_2017}. 
Node kemudian menjalankan bukti kerja ringan (\emph{lightweight proof-of-work}) dengan mencari nilai yang menghasilkan hash di bawah ambang tertentu untuk mencegah spam dan serangan mudah terhadap jaringan, lalu menyiarkan transaksi baru tersebut ke node lain \parencite{TheTangle_Popov_2017,IOTA_Tangle2_2022}. 
Seiring bertambahnya transaksi yang secara langsung atau tidak langsung mereferensikan sebuah transaksi, bobot kumulatif dan tingkat konfirmasinya meningkat hingga transaksi tersebut dianggap dikonfirmasi dengan probabilitas yang cukup tinggi \parencite{TheTangle_Popov_2017}. 

IOTA 2.0 memperkenalkan sejumlah peningkatan terhadap mekanisme konsensus Tangle dengan tujuan menghilangkan peran \emph{coordinator} dan meningkatkan desentralisasi \parencite{IOTA_Tangle2_2022}. 
Dalam arsitektur ini, perangkat lunak node dilengkapi dengan komponen seperti modul anti-spam berbasis \emph{mana}, mekanisme pemesanan transaksi, dan protokol voting yang memungkinkan node mencapai konsensus mengenai tumpukan transaksi yang saling bertentangan secara terdistribusi \parencite{IOTA_Tangle2_2022,FromTangleToRebased_2024}. 
Langkah-langkah ini bertujuan mempertahankan sifat \emph{feeless} dan skalabilitas tinggi, sambil mengurangi ketergantungan pada entitas pusat dalam menentukan validitas final transaksi \parencite{IOTA_Tangle2_2022}. 

Studi terbaru kemudian membandingkan Tangle 2.0 dengan arsitektur Rebased yang diperkenalkan IOTA, yang menggantikan model DAG feeless murni dengan ledger objek yang diprogram menggunakan mesin virtual dan mekanisme konsensus \emph{delegated proof-of-stake} (DPoS) \parencite{FromTangleToRebased_2024}. 
Analisis tersebut menunjukkan bahwa Rebased meningkatkan programabilitas dan kinerja untuk kasus penggunaan tertentu melalui eksekusi kontrak pintar yang lebih kaya, namun memperkenalkan kembali biaya transaksi dan pola desentralisasi yang berbeda dibandingkan Tangle 2.0 \parencite{FromTangleToRebased_2024}. 
Bagi skenario IoT dan kesehatan yang sangat sensitif terhadap biaya dan overhead, pemilihan versi protokol IOTA dan konfigurasi jaringan menjadi keputusan desain penting yang harus diselaraskan dengan kebutuhan arsitektur sistem seperti DecMed \parencite{Elevating_2025_iota_ehealth,FromTangleToRebased_2024}. 

Pada konteks data kesehatan, sejumlah karya menunjukkan bahwa Tangle dapat dimanfaatkan untuk merekam hash dan metadata data klinis atau PGHD sebagai transaksi, sementara data lengkapnya disimpan di sistem RME atau penyimpanan terdistribusi lainnya \parencite{Elevating_2025_iota_ehealth,Exploring_IOTA_HealthData_2022}. 
Dengan cara ini, setiap pembaruan atau akses terhadap data kesehatan dapat dikaitkan dengan transaksi IOTA yang memuat identitas kriptografis, timestamp, dan referensi data, sehingga memungkinkan verifikasi integritas dan penyusunan jejak audit tanpa memindahkan seluruh muatan data ke ledger \parencite{Elevating_2025_iota_ehealth,Exploring_IOTA_HealthData_2022}. 
Integrasi IOTA dengan mekanisme kontrol akses berbasis \emph{capabilities} juga memungkinkan pendelegasian hak akses yang lebih halus, di mana token kemampuan didistribusikan dan diaudit melalui Tangle untuk mengatur siapa yang dapat membaca atau menulis data tertentu \parencite{IOTA_AccessControl_2022}. 

\subsubsection{Karakteristik IOTA yang Relevan untuk Data Kesehatan}

Sejumlah karakteristik IOTA menjadikannya menarik untuk pemanfaatan pada domain kesehatan, khususnya untuk aliran data berskala kecil namun berfrekuensi tinggi seperti PGHD dan data sensor \parencite{Elevating_2025_iota_ehealth}. 
Pertama, sifat \emph{feeless} menghilangkan biaya transaksi per pesan, sehingga lebih realistis untuk skenario di mana perangkat menghasilkan banyak transaksi berukuran kecil setiap hari tanpa membebani pasien atau penyedia layanan dengan biaya gas atau \emph{transaction fee} \parencite{Elevating_2025_iota_ehealth}. 
Kedua, model konsensus terdistribusi memungkinkan pemrosesan transaksi secara paralel, yang dapat meningkatkan skalabilitas dibandingkan mekanisme berbasis blok yang memproses transaksi secara batch \parencite{FromTangleToRebased_2024}. 

Dari sisi keamanan dan integritas, Tangle tetap memanfaatkan kriptografi kunci publik untuk penandatanganan transaksi dan pengalamatan, sehingga setiap transaksi data kesehatan yang dicatat dapat ditelusuri asal-usul dan keabsahannya \parencite{Elevating_2025_iota_ehealth}. 
Beberapa studi mengusulkan model penyimpanan di mana data medis atau PGHD disimpan di penyimpanan terdistribusi atau sistem RME, sementara hash dan metadata integritasnya disimpan pada IOTA untuk menyediakan bukti terhadap perubahan yang tidak sah \parencite{Elevating_2025_iota_ehealth}. 
Pendekatan ini memungkinkan pemisahan jelas antara lapisan penyimpanan data klinis dan lapisan ledger, serta mengurangi beban kapasitas Tangle yang pada dasarnya lebih efisien untuk payload yang relatif kecil \parencite{Elevating_2025_iota_ehealth,FromTangleToRebased_2024}. 

\subsubsection{Tantangan dan Pertimbangan Teknis IOTA dalam E-Health}

Walaupun menawarkan sejumlah keunggulan, penerapan IOTA pada domain kesehatan juga disertai tantangan \parencite{Elevating_2025_iota_ehealth}. 
Pertama, sebagai teknologi yang relatif baru dibandingkan blockchain tradisional, ekosistem perangkat, pustaka, serta praktik terbaik untuk integrasi dengan sistem RME dan perangkat medis masih dalam tahap berkembang, sehingga memerlukan perancangan arsitektur yang hati-hati \parencite{Elevating_2025_iota_ehealth}. 
Kedua, aspek desentralisasi dan keamanan masih terus dimatangkan melalui evolusi protokol seperti Tangle 2.0 dan Rebased, yang berarti implementasi di domain kritikal seperti kesehatan perlu mempertimbangkan versi dan konfigurasi jaringan IOTA yang stabil \parencite{FromTangleToRebased_2024}. 

Selain itu, desain skema kontrol akses di atas IOTA perlu memadukan mekanisme kriptografi (misalnya \emph{public key infrastructure} dan token kemampuan) dengan kebijakan klinis serta regulasi yang berlaku \parencite{Elevating_2025_iota_ehealth,IOTA_AccessControl_2022}. 
Sejumlah karya mengusulkan integrasi konsep \emph{capability-based access control} (CapBAC) dengan IOTA, di mana token kemampuan dan rekaman kebijakan disimpan atau direferensikan melalui Tangle, namun desain ini menuntut pengelolaan kunci dan identitas yang baik agar tidak menimbulkan titik lemah baru \parencite{IOTA_AccessControl_2022}. 
Oleh karena itu, meskipun IOTA menyediakan basis DLT yang ringan dan \emph{feeless}, lapisan aplikasi di atasnya tetap harus dirancang dengan memperhatikan prinsip keamanan, privasi, dan kepatuhan regulasi kesehatan \parencite{Elevating_2025_iota_ehealth}. 

\subsubsection{Perbandingan Teknis Blockchain Konvensional dan IOTA}

Untuk merangkum perbedaan utama antara blockchain konvensional dan IOTA dalam konteks pengelolaan data kesehatan, Tabel~\ref{tab:blockchain-iota-comparison} menyajikan perbandingan dimensi teknis yang relevan.

\begin{table}[h]
  \centering
  \caption{Perbandingan teknis blockchain konvensional dan IOTA dalam konteks data kesehatan}
  \label{tab:blockchain-iota-comparison}
  \begin{tabular}{p{3cm}p{5.5cm}p{5.5cm}}
    \toprule
    Aspek & Blockchain konvensional & IOTA / Tangle \\
    \midrule
    Struktur data & Rantai blok linier, setiap blok berisi sejumlah transaksi dan hash blok sebelumnya \parencite{Mayer_2020_ehr_blockchain_slr, Yang_2021_blockchain_architecture_ehr} & DAG (\emph{Tangle}) di mana setiap transaksi mereferensikan dan memvalidasi beberapa transaksi sebelumnya \parencite{TheTangle_Popov_2017,FromTangleToRebased_2024}. \\
    \midrule
    Mekanisme konsensus & Umumnya berbasis PoW pada blockchain publik, atau PBFT/varian lain pada blockchain berizin; throughput terbatas dan konsumsi energi relatif tinggi untuk PoW \parencite{Mayer_2020_ehr_blockchain_slr,Ettaloui_2024_blockchain_ehr_slr} & Konsensus terdistribusi melalui pemilihan tips dan validasi transaksi oleh setiap pengirim; IOTA 2.0 menambahkan mekanisme voting dan anti-spam berbasis \emph{mana} untuk mencapai konsensus tanpa penambang khusus \parencite{IOTA_Tangle2_2022,FromTangleToRebased_2024}. \\
    \midrule
    Biaya transaksi & Umumnya terdapat biaya transaksi, terutama pada jaringan publik, sehingga kurang ideal untuk banyak transaksi kecil berfrekuensi tinggi \parencite{Ettaloui_2024_blockchain_ehr_slr, Rawat2025_Blockchain_Secure_EHR} & Dirancang sebagai \emph{feeless} pada Tangle generasi awal dan IOTA 2.0, sehingga lebih cocok untuk skenario IoT dan PGHD; arsitektur Rebased memperkenalkan kembali biaya transaksi dengan trade-off peningkatan programabilitas \parencite{Elevating_2025_iota_ehealth,FromTangleToRebased_2024}. \\
    \midrule
    Skalabilitas & Throughput dan latensi bergantung pada ukuran blok, frekuensi pembuatan blok, dan mekanisme konsensus; perlu optimasi tambahan (sharding, off-chain) untuk skala besar \parencite{Mayer_2020_ehr_blockchain_slr} & Potensi skalabilitas meningkat seiring bertambahnya jumlah transaksi karena validasi dilakukan paralel oleh pengirim transaksi; desain protokol menargetkan lingkungan IoT bersumber daya terbatas \parencite{TheTangle_Popov_2017,IOTA_Tangle2_2022}. \\
    \midrule
    Penyimpanan data kesehatan & Data klinis besar disimpan off-chain (database, storage terdistribusi); ledger menyimpan hash, pointer, dan metadata untuk integritas dan audit \parencite{Ettaloui_2024_blockchain_ehr_slr,Yang_2021_blockchain_architecture_ehr} & Pola serupa: data kesehatan/PGHD disimpan off-ledger, sementara Tangle menyimpan referensi, hash, dan metadata yang mengikat transaksi dengan data tersebut \parencite{Elevating_2025_iota_ehealth, Exploring_IOTA_HealthData_2022}. \\
    \midrule
    Kesesuaian dengan IoT dan PGHD & Cocok untuk pencatatan peristiwa dengan frekuensi sedang dan kebutuhan keamanan kuat, namun biaya dan latensi dapat menjadi kendala untuk aliran data sensor yang sangat sering \parencite{Ettaloui_2024_blockchain_ehr_slr,Mayer_2020_ehr_blockchain_slr} & Secara khusus dirancang untuk IoT, dengan biaya nol atau sangat rendah per transaksi dan dukungan terhadap perangkat berdaya terbatas, sehingga selaras dengan karakteristik PGHD \parencite{Elevating_2025_iota_ehealth,IOTA_Tangle2_2022}. \\
    \bottomrule
  \end{tabular}
\end{table}

\subsubsection{Ringkasan dan Bridging ke DecMed dan PGHD}

Secara keseluruhan, DLT menyediakan kerangka konseptual bagi tesis ini untuk memahami bagaimana sebuah ledger terdistribusi dapat digunakan sebagai lapisan integritas dan audit terhadap RME, sementara IOTA menawarkan varian DLT yang lebih sesuai dengan karakteristik data IoT dan PGHD berfrekuensi tinggi \parencite{Ettaloui_2024_blockchain_ehr_slr,Elevating_2025_iota_ehealth}. 
Karakteristik IOTA yang \emph{feeless}, terdistribusi, dan berbasis DAG memberikan keuntungan bagi arsitektur sistem yang perlu menangani banyak pesan kecil dari perangkat pasien tanpa menambah overhead biaya transaksi, sekaligus tetap mempertahankan jejak integritas dan aktivitas akses yang dapat diaudit \parencite{Elevating_2025_iota_ehealth,FromTangleToRebased_2024}. 
Dengan demikian, pemilihan IOTA sebagai basis DLT dalam DecMed menjadi relevan untuk mengintegrasikan aliran PGHD ke dalam ekosistem RME terdistribusi yang akan dibahas lebih lanjut pada subbab DecMed serta pembahasan mengenai kebutuhan \emph{data provenance} pada PGHD di bagian berikutnya.